%%%%%%%%%%%%%%%%%%%%%%%%%%%%%%%%%%%%%%%%%%%%%%%%%%%%%%%%%%%%%%%%%%%%%%%%%%
%%%%%% Vatsav Reddy
%%%%%% International Institute of Information Technology, Hyderabad
%%%%%% Presentation to Prof. Praveen Parchuri
%%%%%%%%%%%%%%%%%%%%%%%%%%%%%%%%%%%%%%%%%%%%%%%%%%%%%%%%%%%%%%%%%%%%%%%%%%
\documentclass{beamer}
\hypersetup{pdfpagemode=FullScreen}
\setbeamertemplate{navigation symbols}{}
\usepackage[english]{babel}
\usepackage[latin1]{inputenc}
\usepackage{amsmath,amssymb}
\usepackage[T1]{fontenc}
\usepackage{times}
\usepackage{mathptmx}
\usepackage{graphicx}
\usepackage{booktabs}
\usepackage{hyperref}
\usepackage{xcolor}
\usepackage{algorithm,algorithmic}
\renewcommand{\algorithmicrequire}{\textbf{Input:}}
\renewcommand{\algorithmicensure}{\textbf{Output:}}

\mode<presentation>{
\usetheme{Madrid}
\usecolortheme{whale}
}

\setbeamercovered{transparent}
\setbeamertemplate{bibliography item}[text]
\setbeamertemplate{theorems}[numbered]
\setbeamerfont{title}{size=\large}
\setbeamerfont{date}{size=\tiny}
\setbeamertemplate{itemize items}[ball]
\setbeamertemplate{itemize subitem}[circle]
\setbeamertemplate{itemize subsubitem}[triangle]
\setbeamertemplate{caption}[numbered]

%----------------------------------------------------------------------------------------
%  TITLE PAGE
%----------------------------------------------------------------------------------------

\title[Neural Networks: Foundations]{Foundations of Neural Networks:\\
McCulloch-Pitts, Perceptron \& Gradient Descent}

\author[Vatsav Reddy]{\texorpdfstring{Vatsav Reddy}{Vatsav Reddy}}

\institute[IIIT-H]{
  {\small Presented to}\\
  \medskip
  {\normalsize Prof.\ Praveen Parchuri}\\
  \vspace{0.5cm}
  International Institute of Information Technology, Hyderabad
}

\date{\today}

\begin{document}

\begin{frame}
  \titlepage
\end{frame}

\begin{frame}{Contents}
  \tableofcontents
\end{frame}

\AtBeginSection[]{
  \begin{frame}<beamer>
    \frametitle{Current Section}
    \tableofcontents[currentsection]
  \end{frame}
}

%==========================================================================
\section{McCulloch-Pitts Neuron}
%==========================================================================

\begin{frame}{Biological Motivation}
  \begin{itemize}
    \item \textbf{Dendrite}: receives signals from other neurons
    \item \textbf{Synapse}: point of connection to other neurons
    \item \textbf{Soma}: processes information
    \item \textbf{Axon}: transmits the output of this neuron
  \end{itemize}
  \vspace{0.3cm}
  \begin{block}{Key Fact}
    The brain is a massively parallel interconnected network of $\approx 10^{11}$ neurons,
    arranged in a hierarchy. Each neuron may perform a certain role / respond to a certain stimulus.
  \end{block}
\end{frame}

\begin{frame}{McCulloch-Pitts (MP) Neuron (1943)}
  McCulloch (Neuroscientist) \& Pitts (Logician) proposed a highly simplified computational model.
  \vspace{0.2cm}
  \begin{columns}
    \column{0.55\textwidth}
    \begin{itemize}
      \item Inputs $x_1, x_2, \ldots, x_n \in \{0,1\}$
      \item $g$ aggregates: $\displaystyle g(x) = \sum_{i=1}^{n} x_i$
      \item $f$ takes a binary decision:
        \[
          y = f(g(x)) = \begin{cases} 1 & g(x) \geq \theta \\ 0 & g(x) < \theta \end{cases}
        \]
      \item $\theta$ = \textbf{thresholding parameter} $\Rightarrow$ \emph{Threshold Logic}
    \end{itemize}
    \column{0.4\textwidth}
    \begin{exampleblock}{Inhibitory Input Rule}
      If any $x_i$ is inhibitory, $y = 0$ regardless of threshold.
    \end{exampleblock}
    \vspace{0.2cm}
    \begin{alertblock}{Limitation}
      Only \emph{linearly separable} Boolean functions. XOR cannot be represented!
    \end{alertblock}
  \end{columns}
\end{frame}

%==========================================================================
\section{Perceptron \& Convergence Proof}
%==========================================================================

\begin{frame}{From MP Neuron to Perceptron}
  Four key advances (Rosenblatt):
  \begin{enumerate}
    \item \textbf{Weights} $w_i$ on inputs (inputs are no longer equal)
    \item \textbf{Weight learning mechanism} from data
    \item \textbf{Real-valued inputs} (not limited to Boolean)
    \item Refined \& carefully analyzed by Minsky \& Papert (1969)
  \end{enumerate}
  \vspace{0.3cm}
  \[
    y = \begin{cases} 1 & \sum_i w_i x_i \geq \theta \\ 0 & \sum_i w_i x_i < \theta \end{cases}
  \]
  Setting $x_0 = 1,\; w_0 = -\theta$ (\textbf{bias node}):
  \[
    y = \begin{cases} 1 & \mathbf{w}^T\mathbf{x} \geq 0 \\ 0 & \mathbf{w}^T\mathbf{x} < 0 \end{cases}
  \]
\end{frame}

\begin{frame}{Perceptron Learning Algorithm}
  \begin{algorithm}[H]
  \begin{algorithmic}[1]
    \REQUIRE $P$ = inputs with label 1;\quad $N$ = inputs with label 0
    \ENSURE Weight vector $\mathbf{w}$
    \STATE $P' \leftarrow P \cup N^{-}$ \quad (where $N^{-}$ contains negatives of all points in $N$)
    \STATE Initialize $\mathbf{w}$ randomly
    \WHILE{not converged}
      \STATE Pick random $\mathbf{p} \in P'$;\quad normalize: $\mathbf{p} \leftarrow \mathbf{p}/\|\mathbf{p}\|$
      \IF{$\mathbf{w}^T\mathbf{p} < 0$}
        \STATE $\mathbf{w} \leftarrow \mathbf{w} + \mathbf{p}$
      \ENDIF
    \ENDWHILE
  \end{algorithmic}
  \caption{Perceptron Learning (normalized)}
  \end{algorithm}
  \vspace{0.2cm}
  \begin{itemize}
    \item $\mathbf{w}^T\mathbf{x} = 0$ defines the decision hyperplane; $\mathbf{w}$ is orthogonal to it.
    \item Each correction rotates $\mathbf{w}$ closer (in angle) to the misclassified point.
  \end{itemize}
\end{frame}

\begin{frame}{Convergence Proof --- Setup \& Statement}
  \begin{theorem}[Perceptron Convergence]
    If $P$ and $N$ are finite and \emph{linearly separable}, the perceptron learning
    algorithm makes at most $k_{\max}$ corrections and \textbf{converges} in finite steps.
  \end{theorem}

  \vspace{0.2cm}
  \textbf{Setup:}
  \begin{itemize}
    \item Let $\mathbf{w}^*$ be the (normalized) optimal solution: $\|\mathbf{w}^*\| = 1$.
    \item Since data is linearly separable, $\exists\;\delta > 0$ such that
          $\mathbf{w}^{*T}\mathbf{p} \geq \delta$ for all normalized $\mathbf{p} \in P'$.
    \item Let $\beta_k$ = angle between $\mathbf{w}^*$ and $\mathbf{w}_k$ after $k$ corrections.
    \item Track:
    \[
      \cos\beta_k = \frac{\mathbf{w}^*\cdot\mathbf{w}_k}{\|\mathbf{w}^*\|\,\|\mathbf{w}_k\|}
      = \frac{\mathbf{w}^*\cdot\mathbf{w}_k}{\|\mathbf{w}_k\|}
    \]
  \end{itemize}
\end{frame}

\begin{frame}{Convergence Proof --- Numerator Grows Linearly}
  At step $t$: correction is $\mathbf{w}_{t+1} = \mathbf{w}_t + \mathbf{p}_t$
  (triggered because $\mathbf{w}_t^T\mathbf{p}_t < 0$).

  \vspace{0.2cm}
  \textbf{Effect on numerator $\mathbf{w}^*\cdot\mathbf{w}_{t+1}$:}
  \begin{align*}
    \mathbf{w}^*\cdot\mathbf{w}_{t+1} &= \mathbf{w}^*\cdot(\mathbf{w}_t + \mathbf{p}_t) \\
                                        &= \mathbf{w}^*\cdot\mathbf{w}_t + \underbrace{\mathbf{w}^*\cdot\mathbf{p}_t}_{\geq\;\delta} \\[4pt]
                                        &\geq \mathbf{w}^*\cdot\mathbf{w}_t + \delta
  \end{align*}

  \vspace{0.1cm}
  By induction over $k$ corrections (starting from $\mathbf{w}_0$):
  \[
    \boxed{\mathbf{w}^*\cdot\mathbf{w}_k \;\geq\; \mathbf{w}^*\cdot\mathbf{w}_0 + k\delta}
  \]
  The numerator grows \textbf{at least linearly} in $k$.
\end{frame}

\begin{frame}{Convergence Proof --- Denominator Grows Sub-Linearly}
  \textbf{Effect on $\|\mathbf{w}_{t+1}\|^2$:}
  \begin{align*}
    \|\mathbf{w}_{t+1}\|^2 &= \|\mathbf{w}_t + \mathbf{p}_t\|^2 \\
                             &= \|\mathbf{w}_t\|^2 + 2\underbrace{\mathbf{w}_t^T\mathbf{p}_t}_{<\;0}
                                + \underbrace{\|\mathbf{p}_t\|^2}_{=\;1\;(\text{normalized})} \\[4pt]
                             &\leq \|\mathbf{w}_t\|^2 + 1
  \end{align*}

  By induction after $k$ corrections:
  \[
    \boxed{\|\mathbf{w}_k\|^2 \;\leq\; \|\mathbf{w}_0\|^2 + k}
  \]

  \begin{block}{Observation}
    Denominator $\|\mathbf{w}_k\| \leq \sqrt{\|\mathbf{w}_0\|^2 + k}$, which grows only as
    $\sqrt{k}$ --- \textbf{sub-linearly} in $k$.
  \end{block}
\end{frame}

\begin{frame}{Convergence Proof --- Conclusion}
  Combining both bounds:
  \[
    \cos\beta_k
    = \frac{\mathbf{w}^*\cdot\mathbf{w}_k}{\|\mathbf{w}^*\|\,\|\mathbf{w}_k\|}
    \;\geq\;
    \frac{\mathbf{w}^*\cdot\mathbf{w}_0 + k\delta}{\sqrt{\|\mathbf{w}_0\|^2 + k}}
    \;\xrightarrow{k\to\infty}\; \infty
  \]

  This grows proportionally to $\sqrt{k}\cdot\delta$.

  \begin{alertblock}{Contradiction!}
    $\cos\beta_k \leq 1$ by definition of cosine --- it can never exceed 1.
  \end{alertblock}

  \begin{block}{Therefore}
    $k$ must be \textbf{bounded above} by some finite $k_{\max}$:
    \[
      k \;\leq\; k_{\max} \;=\; \frac{(\mathbf{w}^*\cdot\mathbf{w}_0)^2 + \|\mathbf{w}_0\|^2 \delta^2}{\delta^2}
      \quad (\text{approx.})
    \]
    $\Rightarrow$ Finite corrections $\Rightarrow$ Algorithm \textbf{converges}. \hfill $\square$
  \end{block}
\end{frame}

%==========================================================================
\section{Sigmoid Neuron \& Loss}
%==========================================================================

\begin{frame}{Sigmoid Neuron}
  \begin{columns}
    \column{0.5\textwidth}
    \textbf{Perceptron}: hard threshold, binary output.\\[6pt]
    \textbf{Sigmoid Neuron}: smooth, real-valued output in $(0,1)$:
    \[
      y = \sigma(\mathbf{w}^T\mathbf{x}) = \frac{1}{1+e^{-\mathbf{w}^T\mathbf{x}}}
    \]
    \begin{itemize}
      \item $\mathbf{w}^T\mathbf{x} \to +\infty \;\Rightarrow\; y \to 1$
      \item $\mathbf{w}^T\mathbf{x} \to -\infty \;\Rightarrow\; y \to 0$
      \item $\mathbf{w}^T\mathbf{x} = 0 \;\Rightarrow\; y = \tfrac{1}{2}$
    \end{itemize}
    \column{0.45\textwidth}
    \begin{block}{Key Property}
      \[
        \sigma'(x) = \sigma(x)(1 - \sigma(x))
      \]
      Output as \emph{probability} enables gradient-based learning --- no sharp jump around threshold $(-w_0)$.
    \end{block}
    \begin{exampleblock}{Loss Function}
      \[
        \mathcal{L}(\mathbf{w},b)
        = \frac{1}{N}\sum_{i=1}^{N}\bigl(y_i - \sigma(x_i)\bigr)^2
      \]
    \end{exampleblock}
  \end{columns}
\end{frame}

%==========================================================================
\section{Taylor Series \& Gradient Descent Proofs}
%==========================================================================

\begin{frame}{Taylor Series --- The Foundation}
  \textbf{Scalar form:}
  \[
    f(x+h) = f(x) + hf'(x) + \frac{h^2}{2}f''(x) + \cdots
  \]
  Rewriting with $y = x+h$:
  \[
    f(y) = f(x) + (y-x)f'(x) + \frac{(y-x)^2}{2}f''(x) + \cdots
  \]
  \textbf{Multivariate form} (gradient $\nabla f$, Hessian $H$):
  \[
    f(y) = f(x) + (y-x)^T\nabla f(x) + \frac{1}{2}(y-x)^T H(x)(y-x) + \cdots
  \]
  Setting $\mathbf{w}_{k+1} = \mathbf{w}_k + s$:
  \[
    \boxed{f(\mathbf{w}_{k+1}) = f(\mathbf{w}_k) + s^T\nabla f(\mathbf{w}_k)
    + \frac{1}{2}s^T H s + \cdots}
  \]
  \small{(Here $\nabla f$ and $H$ are evaluated at $\mathbf{w}_k$.)}
\end{frame}

\begin{frame}{GD via Taylor Series}
  Apply to loss $J(\mathbf{w})$ with step $s = \mathbf{w}_{k+1} - \mathbf{w}_k$:
  \[
    J(\mathbf{w}_{k+1})
    = J(\mathbf{w}_k) + (\mathbf{w}_{k+1}-\mathbf{w}_k)^T\nabla J(\mathbf{w}_k)
    + \frac{1}{2}(\mathbf{w}_{k+1}-\mathbf{w}_k)^T H (\mathbf{w}_{k+1}-\mathbf{w}_k) + \cdots
  \]
  Compactly:
  \[
    J(\mathbf{w}_{k+1}) = J(\mathbf{w}_k) + s^T\nabla J + \frac{1}{2}s^T H s + \cdots
  \]
  \begin{block}{Gradient Descent Update}
    \[
      s \;=\; -\eta\nabla J(\mathbf{w}_k)
      \quad\Longleftrightarrow\quad
      \mathbf{w}_{k+1} = \mathbf{w}_k - \eta\,\nabla J(\mathbf{w}_k)
    \]
  \end{block}
\end{frame}

\begin{frame}{Why Does GD Work? --- Non-Increasing Updates}
  Use the \textbf{first-order} Taylor approximation and substitute $s = -\eta\nabla J$:
  \begin{align*}
    J(\mathbf{w}_{k+1}) &\approx J(\mathbf{w}_k) + s^T\nabla J(\mathbf{w}_k) \\[4pt]
                         &= J(\mathbf{w}_k) + (-\eta\nabla J)^T \nabla J \\[4pt]
                         &= J(\mathbf{w}_k) - \eta\underbrace{\|\nabla J(\mathbf{w}_k)\|^2}_{\geq\;0}
  \end{align*}
  \begin{block}{Conclusion}
    For any $\eta > 0$:
    \[
      J(\mathbf{w}_{k+1}) \leq J(\mathbf{w}_k)
      \quad\Rightarrow\quad \text{\textbf{Non-increasing} at every step.}
    \]
    When $\nabla J = \mathbf{0}$: no change in objective $\Rightarrow$ iterations stop at a stationary point.
  \end{block}
\end{frame}

\begin{frame}{Why Steepest Descent? --- Directional Derivative}
  From the first-order Taylor expansion:
  \[
    J(\mathbf{x}+\Delta\mathbf{x}) \approx J(\mathbf{x}) + \Delta\mathbf{x}^T\nabla J(\mathbf{x})
  \]
  The \textbf{directional derivative} along unit vector $\mathbf{u}$ is $\mathbf{u}^T\nabla J$.\\[6pt]
  \textbf{Best direction to decrease $J$?}
  \[
    \min_{\mathbf{u},\;\|\mathbf{u}\|=1}\; \mathbf{u}^T\nabla J
    \;=\; \min_{\mathbf{u},\;\|\mathbf{u}\|=1}\;\|\mathbf{u}\|\,\|\nabla J\|\cos\theta
  \]
  Minimized when $\cos\theta = -1$, i.e., $\mathbf{u}$ points \textbf{opposite} to $\nabla J$.
  \begin{block}{Steepest Descent}
    \[
      \mathbf{x}_{\text{new}} = \mathbf{x}_{\text{old}} - \eta\,\nabla J
    \]
    Moving opposite to the gradient gives the \textbf{steepest decrease} in $J$.
  \end{block}
\end{frame}

\begin{frame}{Is There an Optimal Learning Rate?}
  Using the \textbf{second-order} Taylor expansion with $s = -\eta\nabla J$:
  \[
    J(\mathbf{w}_{k+1})
    = J(\mathbf{w}_k) - \eta\nabla J^T\nabla J + \frac{\eta^2}{2}\nabla J^T H\nabla J
  \]
  Differentiate w.r.t.\ $\eta$ and set to zero:
  \[
    \frac{d}{d\eta}J(\mathbf{w}_{k+1})
    = -\|\nabla J\|^2 + \eta\,\nabla J^T H\nabla J = 0
  \]
  \[
    \boxed{\eta^* = \frac{\|\nabla J\|^2}{\nabla J^T H\nabla J}}
  \]
  \begin{alertblock}{Note}
    $\eta^*$ changes at every iteration $k$ since $\nabla J$ and $H$ both depend on $\mathbf{w}_k$.
    Computing $H$ at every step is expensive --- motivates second-order methods.
  \end{alertblock}
\end{frame}

\begin{frame}{Gradient, Hessian \& Jacobian --- Recap}
  \begin{columns}
    \column{0.33\textwidth}
    \begin{block}{Gradient}
      $f:\mathbb{R}^n\to\mathbb{R}$
      \[
        \nabla f = \begin{bmatrix}\frac{\partial f}{\partial x_1}\\\vdots\\\frac{\partial f}{\partial x_n}\end{bmatrix}
      \]
      Scalar output.\\
      Steepest ascent direction.
    \end{block}
    \column{0.33\textwidth}
    \begin{block}{Hessian}
      $f:\mathbb{R}^n\to\mathbb{R}$
      \[
        H_{ij} = \frac{\partial^2 f}{\partial x_i \partial x_j}
      \]
      Curvature matrix.\\
      Tells \emph{how far} to trust gradient.
    \end{block}
    \column{0.33\textwidth}
    \begin{block}{Jacobian}
      $f:\mathbb{R}^n\to\mathbb{R}^m$
      \[
        J_{ij} = \frac{\partial f_i}{\partial x_j}
      \]
      Generalizes gradient to vector-valued functions.
    \end{block}
  \end{columns}
\end{frame}

\begin{frame}{Second-Order Methods (Newton's Method)}
  \textbf{Idea}: find the step $s$ that minimizes the second-order Taylor expansion.\\[4pt]
  Differentiate $J(\mathbf{w}_k) + s^T\nabla J + \frac{1}{2}s^T H s$ w.r.t.\ $s$ and set to zero:
  \[
    \nabla J + H s = 0 \;\Longrightarrow\; s = -H^{-1}\nabla J
  \]
  \[
    \boxed{\mathbf{w}_{k+1} = \mathbf{w}_k - H^{-1}\nabla J(\mathbf{w}_k)}
  \]
  \vspace{0.1cm}
  \begin{columns}
    \column{0.5\textwidth}
    \begin{exampleblock}{Advantages}
      \begin{itemize}
        \item Often \textbf{quadratic convergence}
        \item Second-order approx.\ more accurate
        \item Gradient gives direction; Hessian gives magnitude
      \end{itemize}
    \end{exampleblock}
    \column{0.45\textwidth}
    \begin{alertblock}{Disadvantages}
      \begin{itemize}
        \item $H$ is $d\times d$ --- huge memory
        \item $H^{-1}$ costs $O(d^3)$
        \item Negative eigenvalues (saddle points) cause ascent steps
        \item Fix: regularize as $(H+\alpha I)^{-1}\nabla J$
      \end{itemize}
    \end{alertblock}
  \end{columns}
\end{frame}

%==========================================================================
\section{Summary}
%==========================================================================

\begin{frame}{Summary}
  \begin{columns}
    \column{0.48\textwidth}
    \begin{block}{Neural Network Foundations}
      \begin{itemize}
        \item \textbf{MP Neuron}: threshold logic on Boolean inputs; only linearly separable functions
        \item \textbf{Perceptron}: weighted inputs + bias + learning; converges if data is linearly separable (proved via $\cos\beta$)
        \item \textbf{Sigmoid}: smooth probabilistic output enabling gradient flow
      \end{itemize}
    \end{block}
    \column{0.48\textwidth}
    \begin{block}{Optimization Foundations}
      \begin{itemize}
        \item \textbf{Taylor Series}: base for all GD analysis
        \item \textbf{GD}: non-increasing updates; steepest descent is optimal 1st-order direction
        \item \textbf{Optimal $\eta$}: from 2nd-order Taylor
        \item \textbf{Newton's}: faster but costly $H^{-1}$; regularize for saddle points
      \end{itemize}
    \end{block}
  \end{columns}
  \vspace{0.3cm}
  \begin{center}
    \large\textit{``Gradient tells direction; Hessian tells how much to trust it.''}
  \end{center}
\end{frame}

\begin{frame}
  \Huge{\centerline{Thank You}}
  \vspace{0.5cm}
  \large{\centerline{Questions?}}
\end{frame}

\end{document}